\documentclass[11pt,a4paper]{article}

% --- Encoding ---
\usepackage[utf8]{inputenc}
\usepackage[T1]{fontenc}
\usepackage[spanish,es-noshorthands]{babel}
\usepackage{lmodern}

% --- Layout ---
\usepackage[margin=2.5cm]{geometry}
\usepackage{parskip}
\setlength{\parindent}{0pt}
\setlength{\parskip}{0.5em plus 0.2em minus 0.1em}

% --- Tables and graphics ---
\usepackage{booktabs}
\usepackage{array}
\usepackage{enumitem}
\usepackage{xcolor}
\usepackage{graphicx}
\usepackage{tikz}
\usetikzlibrary{positioning,shapes.geometric,arrows.meta,fit,backgrounds,calc,decorations.pathreplacing}
\usepackage{caption}
\usepackage{subcaption}
\usepackage{amsmath,amssymb}

% --- Code listings ---
\usepackage{listings}
\definecolor{codegray}{rgb}{0.96,0.96,0.96}
\definecolor{codeborder}{rgb}{0.85,0.85,0.85}
\definecolor{codecomment}{rgb}{0.30,0.55,0.30}
\definecolor{codekeyword}{rgb}{0.10,0.30,0.70}
\definecolor{codestring}{rgb}{0.65,0.15,0.10}

\lstdefinestyle{cstyle}{
    language=C,
    backgroundcolor=\color{codegray},
    basicstyle=\ttfamily\footnotesize,
    breaklines=true,
    frame=single,
    rulecolor=\color{codeborder},
    showstringspaces=false,
    keywordstyle=\color{codekeyword}\bfseries,
    commentstyle=\color{codecomment}\itshape,
    stringstyle=\color{codestring},
    aboveskip=1em,
    belowskip=1em,
}

\lstdefinestyle{shellstyle}{
    language=bash,
    backgroundcolor=\color{codegray},
    basicstyle=\ttfamily\small,
    breaklines=true,
    frame=single,
    rulecolor=\color{codeborder},
    showstringspaces=false,
    keywordstyle=\color{codekeyword}\bfseries,
    commentstyle=\color{codecomment}\itshape,
    stringstyle=\color{codestring},
    aboveskip=1em,
    belowskip=1em,
    columns=fullflexible,
}

\lstdefinestyle{textstyle}{
    backgroundcolor=\color{codegray},
    basicstyle=\ttfamily\footnotesize,
    breaklines=true,
    frame=single,
    rulecolor=\color{codeborder},
    showstringspaces=false,
    aboveskip=1em,
    belowskip=1em,
    columns=fullflexible,
}

\lstset{style=textstyle}

% --- Definition / Idea / Important boxes ---
\usepackage[most]{tcolorbox}
\newtcolorbox{definicion}[1][]{
    colback=blue!5!white,
    colframe=blue!50!black,
    fonttitle=\bfseries,
    title=Definición,
    sharp corners=southwest,
    rounded corners=northwest,
    arc=2pt,
    boxrule=0.6pt,
    left=8pt,right=8pt,top=4pt,bottom=4pt,
    #1
}
\newtcolorbox{idea}[1][]{
    colback=orange!5!white,
    colframe=orange!60!black,
    fonttitle=\bfseries,
    title=Idea intuitiva,
    sharp corners=southwest,
    rounded corners=northwest,
    arc=2pt,
    boxrule=0.6pt,
    left=8pt,right=8pt,top=4pt,bottom=4pt,
    #1
}
\newtcolorbox{importante}[1][]{
    colback=red!5!white,
    colframe=red!50!black,
    fonttitle=\bfseries,
    title=Importante,
    sharp corners=southwest,
    rounded corners=northwest,
    arc=2pt,
    boxrule=0.6pt,
    left=8pt,right=8pt,top=4pt,bottom=4pt,
    #1
}

% --- Hyperlinks ---
\usepackage[colorlinks=true,
            linkcolor=blue!60!black,
            urlcolor=blue!50!black,
            citecolor=blue!60!black]{hyperref}

% --- Memory palette ---
\definecolor{mem1}{HTML}{4C72B0}
\definecolor{mem2}{HTML}{DD8452}
\definecolor{mem3}{HTML}{55A868}
\definecolor{mem4}{HTML}{8172B2}
\definecolor{memused}{HTML}{C44E52}
\definecolor{memfree}{HTML}{8C8C8C}
\definecolor{memshared}{HTML}{4FC3C7}
\definecolor{memfile}{HTML}{F4A261}

% --- Title block ---
\title{
    \vspace{-2.5em}
    \textbf{Práctico 5} \\[0.3em]
    \large Gestión de la memoria y memoria virtual
}
\author{
    Tomás Rodeghiero \\
    \small Sistemas Operativos --- Universidad Nacional de Río Cuarto
}
\date{2026}

\begin{document}

\maketitle

\begin{abstract}
\noindent
Este documento presenta la resolución integral del Práctico 5 de la
materia Sistemas Operativos (UNRC, 2026), centrado en la
\emph{gestión de la memoria} y la \emph{memoria virtual}. Se resuelven
nueve ejercicios que cubren: inspección del consumo de memoria con
\texttt{top}, \texttt{free}, \texttt{smem} y \texttt{swapon};
\emph{buddy system}; análisis del layout lógico de un proceso C;
arquitecturas con paginado (tabla de páginas, formato de \emph{pte},
páginas grandes vs.\ multinivel); memoria compartida entre procesos;
mapeo de archivos en memoria (\texttt{mmap}); y la API de \emph{shared
memory} de UNIX (\texttt{shmget}, \texttt{shmat}, \texttt{shmdt},
\texttt{shmctl}). Cada ejercicio se justifica con los capítulos
\emph{Gestión de la memoria} y \emph{Memoria virtual} de las Notas del
curso (Marcelo Arroyo) y con los resúmenes teóricos
\texttt{resumen-teorico-gestion-memoria} y
\texttt{resumen-teorico-memoria-virtual} disponibles en
\texttt{materiales/}.
\end{abstract}

\tableofcontents
\bigskip

% =====================================================================
\section{Introducción}
% =====================================================================

La gestión de la memoria es una de las funciones más importantes que
realiza el sistema operativo. Sus dos responsabilidades inseparables
son la \textbf{asignación} (decidir qué regiones de la RAM le tocan a
qué procesos y al kernel, y cómo se liberan cuando ya no se usan) y la
\textbf{protección} (garantizar que un proceso se ejecute confinado a
sus regiones, sin poder invadir las de otros procesos ni las del
kernel).

A su vez, la \textbf{memoria virtual} es la técnica que permite
ejecutar programas cuya memoria total excede la RAM disponible,
manteniendo en RAM sólo las partes en uso en cada momento y dejando el
resto en disco. Se apoya en el esquema de \emph{paginado} y en bits de
control de la \texttt{pte} (\texttt{V}, \texttt{A}, \texttt{D}) para
implementar \emph{swapping}/\emph{demand paging}, \texttt{mmap},
crecimiento dinámico del stack, memoria compartida entre procesos y
\emph{copy on write}.

Este documento aplica esos conceptos a problemas concretos: desde
inspeccionar el consumo real del sistema hasta diseñar tablas de
páginas para una arquitectura dada, pasando por la simulación de un
\emph{buddy allocator} y el análisis empírico de \texttt{mmap} y la
memoria compartida en Linux.

% =====================================================================
\section{Ejercicio 1 --- Memoria y swap usados por los procesos}
% =====================================================================

\textbf{Enunciado.} Determinar los espacios de memoria (incluyendo
swap) que están usando los procesos del sistema. Ayuda: ver los
comandos \texttt{top}, \texttt{free}, \texttt{smem} y \texttt{swapon}.

\subsection{Resumen global con \texttt{free}}

\texttt{free} presenta una visión \emph{global} del sistema en una
sola línea: cuánta RAM y cuánto \emph{swap} hay, cuánto está usado,
libre, compartido y cuánto se ocupa en \emph{buffers/cache}.

\begin{lstlisting}[style=shellstyle]
$ free -h
              total      usado      libre      compart.    buff/cache  disponible
Mem:           7.6Gi      3.1Gi      980Mi       340Mi       3.6Gi      3.8Gi
Swap:          2.0Gi      512Mi      1.5Gi
\end{lstlisting}

Interpretación de columnas:
\begin{itemize}[leftmargin=*,nosep]
    \item \textbf{total}: capacidad total de RAM (o swap) reconocida
        por el kernel.
    \item \textbf{usado}: RAM efectivamente asignada a procesos,
        kernel y reservada.
    \item \textbf{libre}: páginas libres (con bit \texttt{V} apagado
        en ningún \emph{pte} que las refiera).
    \item \textbf{compart.}: páginas de \emph{shared memory}
        (\texttt{tmpfs}, IPC, \texttt{mmap MAP\_SHARED}).
    \item \textbf{buff/cache}: páginas usadas por el cache de disco
        (\emph{page cache}). Son liberables: el kernel las descarta
        si la RAM escasea.
    \item \textbf{disponible}: estimación de cuánta memoria se puede
        asignar sin recurrir a swap (esencialmente
        \texttt{libre + buff/cache} reclamable).
\end{itemize}

\subsection{Dispositivos de swap con \texttt{swapon}}

\texttt{swapon --show} lista los dispositivos o archivos configurados
como swap. Cada entrada indica tipo (\emph{partition} o \emph{file}),
tamaño y cuánto se está usando:

\begin{lstlisting}[style=shellstyle]
$ swapon --show
NAME      TYPE        SIZE     USED    PRIO
/dev/sda3 partition   2.0G     512M    -2
\end{lstlisting}

Una \emph{partición} con formato de tipo \emph{swap} es más rápida que
un \emph{archivo} swap, porque evita la indirección del sistema de
archivos.

\subsection{Vista por proceso con \texttt{top}}

\texttt{top} muestra los procesos en tiempo real. Las columnas
relevantes para memoria son:

\begin{center}
\small
\begin{tabular}{@{}lp{11cm}@{}}
\toprule
Columna & Significado \\
\midrule
\texttt{VIRT} & Tamaño total del espacio de direcciones lógicas del proceso (texto, datos, heap, stack, \texttt{mmap}, librerías compartidas, etc.). \\
\texttt{RES}  & \emph{Resident Set Size}: páginas que efectivamente están en RAM. \\
\texttt{SHR}  & Páginas residentes que son \emph{compartidas} con otros procesos (típicamente librerías \texttt{.so}, \texttt{mmap MAP\_SHARED}). \\
\texttt{\%MEM} & \texttt{RES} dividido por la RAM total. \\
\texttt{SWAP} & Páginas del proceso que actualmente residen en el área de swap (sólo si se habilita la columna con \texttt{f} en \texttt{top}). \\
\bottomrule
\end{tabular}
\end{center}

\begin{importante}
\texttt{VIRT} \emph{no} representa memoria realmente consumida: incluye
\emph{mappings} (por ejemplo, librerías compartidas) que pueden no
estar en RAM. Sumar \texttt{VIRT} de todos los procesos puede dar
varias veces la RAM física y no significa que el sistema esté
\emph{out of memory}.
\end{importante}

\subsection{Vista detallada con \texttt{smem}}

\texttt{smem} es la herramienta más precisa para conocer el consumo
\emph{real} de cada proceso, porque introduce las métricas
\textbf{USS} y \textbf{PSS}:

\begin{itemize}[leftmargin=*]
    \item \textbf{USS} (\emph{Unique Set Size}): memoria que sólo
        usa el proceso. Es lo que se liberaría si el proceso terminara.
    \item \textbf{PSS} (\emph{Proportional Set Size}): RAM única más
        las páginas compartidas \emph{prorrateadas} entre los procesos
        que las comparten. La suma de PSS de todos los procesos da el
        consumo real de memoria de usuario.
    \item \textbf{RSS} (\emph{Resident Set Size}): RAM total mapeada
        por el proceso (cuenta cada página compartida en cada proceso
        que la usa).
    \item \textbf{SWAP}: páginas del proceso en el área de swap.
\end{itemize}

\begin{lstlisting}[style=shellstyle]
$ smem -tk -c "pid user command swap uss pss rss"
  PID User  Command                    Swap     USS     PSS     RSS
 1234 tomi  /usr/bin/firefox            12M   320M    430M    920M
 4567 tomi  /usr/lib/code/code-helper    8M    90M    140M    400M
 ...
-----------------------------------------------------------------
                                       128M  1.1G   1.6G    3.4G
\end{lstlisting}

\begin{idea}
\texttt{RSS} sobreestima el consumo porque cuenta varias veces la
misma página compartida; \texttt{USS} subestima porque ignora las
compartidas; \texttt{PSS} es el mejor compromiso. Cuando se pregunta
``\textquestiondown{}cuánta RAM consume este proceso de verdad?'', la respuesta más
honesta es PSS.
\end{idea}

\subsection{Información cruda por proceso}

El kernel expone los datos por proceso en \texttt{/proc/<pid>/status}
y \texttt{/proc/<pid>/smaps}:

\begin{lstlisting}[style=shellstyle]
$ grep -E "VmSize|VmRSS|VmSwap|VmData|VmStk|VmExe|VmLib" /proc/$$/status
VmSize:    23012 kB     # tamano virtual total
VmRSS:      4140 kB     # resident set size
VmData:     1872 kB     # segmento de datos
VmStk:       132 kB     # stack
VmExe:       872 kB     # codigo (text)
VmLib:      3568 kB     # librerias compartidas
VmSwap:      256 kB     # paginas en swap
\end{lstlisting}

El archivo \texttt{/proc/<pid>/smaps} lista cada VMA (región) por
separado con campos como \texttt{Size}, \texttt{Rss}, \texttt{Pss},
\texttt{Swap}, \texttt{Shared\_Clean} y \texttt{Shared\_Dirty}, entre
otros.

\subsection*{Conclusión}

\begin{itemize}[leftmargin=*,nosep]
    \item \texttt{free} y \texttt{swapon} dan el panorama global del
        sistema.
    \item \texttt{top} ordena los procesos por consumo aparente
        (\texttt{VIRT}/\texttt{RES}) pero no diferencia memoria
        compartida.
    \item \texttt{smem} agrega \texttt{PSS}/\texttt{USS} y es lo más
        adecuado para responder por proceso.
    \item \texttt{/proc/<pid>/status} y \texttt{smaps} son la fuente
        cruda que usan todas las herramientas anteriores.
\end{itemize}

% =====================================================================
\section{Ejercicio 2 --- Evolución de un heap con \emph{buddy system}}
% =====================================================================

\textbf{Enunciado.} Mostrar la evolución de un heap de tamaño $1024$
($2^{10}$) gestionado con un \emph{buddy system} con bloques de tamaño
entre $2^5$ y $2^{10}$ ante la secuencia
\texttt{ptr1=malloc(20)},
\texttt{ptr2=malloc(50)},
\texttt{ptr3=malloc(1010)},
\texttt{free(ptr1)} y \texttt{malloc(64)}.

\paragraph{Sobre el tamaño del heap.} El enunciado original menciona
``$1$\,MB'', pero el rango de bloques permitidos es $2^5$ ($32$ bytes)
a $2^{10}$ ($1024$ bytes). La interpretación más coherente con esa
restricción es que el heap tiene tamaño $2^{10} = 1024$ bytes (es decir,
un único bloque máximo). Así, todas las operaciones se pueden expresar
en términos de los tamaños permitidos y la subdivisión del bloque
maestro de $1024$ bytes.

\subsection{Recordatorio del algoritmo}

\begin{itemize}[leftmargin=*]
    \item Inicialmente, el heap se particiona en bloques del tamaño
        máximo $2^h$ ($1024$ en este caso).
    \item Ante un pedido de tamaño $n$: se busca el menor bloque
        libre $b$ con $b.\textit{size} \geq 2^{\lceil \log_2 n \rceil}$.
        Si $b.\textit{size} \geq 2n$ y $b.\textit{size} > 2^l$, se
        divide en dos mitades (\emph{buddies}) y se repite el paso.
        Finalmente se asigna $b$.
    \item Ante una liberación: si el \emph{buddy} contiguo del mismo
        tamaño también está libre, se fusionan en un bloque del doble
        de tamaño. La fusión se repite mientras se pueda y sin pasar
        de $2^h$.
\end{itemize}

\subsection{Operación \texttt{ptr1 = malloc(20)}}

$20$ cabe en el menor bloque permitido: $2^5 = 32$ bytes. El único
bloque libre es el de $1024$, así que hay que dividir sucesivamente
hasta llegar a $32$:

\begin{center}
\small
\begin{tabular}{@{}lll@{}}
\toprule
Paso & Acción & Estado del heap (tamaños) \\
\midrule
0 & inicio                        & $[\,1024_L\,]$ \\
1 & divido 1024                   & $[\,512_L \mid 512_L\,]$ \\
2 & divido el primer 512          & $[\,256_L \mid 256_L \mid 512_L\,]$ \\
3 & divido el primer 256          & $[\,128_L \mid 128_L \mid 256_L \mid 512_L\,]$ \\
4 & divido el primer 128          & $[\,64_L \mid 64_L \mid 128_L \mid 256_L \mid 512_L\,]$ \\
5 & divido el primer 64           & $[\,32_L \mid 32_L \mid 64_L \mid 128_L \mid 256_L \mid 512_L\,]$ \\
6 & asigno el primer 32 a ptr1    & $[\,\mathbf{32_U} \mid 32_L \mid 64_L \mid 128_L \mid 256_L \mid 512_L\,]$ \\
\bottomrule
\end{tabular}
\end{center}

\begin{center}
\begin{tikzpicture}[
    every node/.style={font=\ttfamily\scriptsize},
    scale=0.85,transform shape,block/.style={draw,minimum height=0.7cm,align=center}
]
\node[block,fill=memused!40,minimum width=0.5cm] (b0) at (0.25,0) {32};
\node[block,fill=memfree!20,minimum width=0.5cm,right=0pt of b0] (b1) {32};
\node[block,fill=memfree!20,minimum width=1.0cm,right=0pt of b1] (b2) {64};
\node[block,fill=memfree!20,minimum width=2.0cm,right=0pt of b2] (b3) {128};
\node[block,fill=memfree!20,minimum width=4.0cm,right=0pt of b3] (b4) {256};
\node[block,fill=memfree!20,minimum width=8.0cm,right=0pt of b4] (b5) {512};
\node[font=\scriptsize] at ($(b0.south)+(0,-0.3)$) {0};
\node[font=\scriptsize] at ($(b5.south east)+(0,-0.3)$) {1024};
\node[font=\bfseries] at ($(b0.north)+(0,0.25)$) {ptr1};
\end{tikzpicture}
\end{center}

\texttt{ptr1} apunta al offset $0$.

\subsection{Operación \texttt{ptr2 = malloc(50)}}

$50$ cabe en $2^6 = 64$ bytes. Hay un bloque libre de exactamente $64$
(offset $64$), así que se asigna directamente sin más divisiones.

\begin{center}
\begin{tikzpicture}[
    every node/.style={font=\ttfamily\scriptsize},
    scale=0.85,transform shape,block/.style={draw,minimum height=0.7cm,align=center}
]
\node[block,fill=memused!40,minimum width=0.5cm] (b0) at (0.25,0) {32};
\node[block,fill=memfree!20,minimum width=0.5cm,right=0pt of b0] (b1) {32};
\node[block,fill=memused!40,minimum width=1.0cm,right=0pt of b1] (b2) {64};
\node[block,fill=memfree!20,minimum width=2.0cm,right=0pt of b2] (b3) {128};
\node[block,fill=memfree!20,minimum width=4.0cm,right=0pt of b3] (b4) {256};
\node[block,fill=memfree!20,minimum width=8.0cm,right=0pt of b4] (b5) {512};
\node[font=\scriptsize] at ($(b0.south)+(0,-0.3)$) {0};
\node[font=\scriptsize] at ($(b5.south east)+(0,-0.3)$) {1024};
\node[font=\bfseries] at ($(b0.north)+(0,0.25)$) {ptr1};
\node[font=\bfseries] at ($(b2.north)+(0,0.25)$) {ptr2};
\end{tikzpicture}
\end{center}

\texttt{ptr2} apunta al offset $64$.

\subsection{Operación \texttt{ptr3 = malloc(1010)}}

$1010$ necesita un bloque de $2^{10} = 1024$ bytes. Pero el bloque
inicial de $1024$ ya fue subdividido y sus pedazos no están
disponibles en forma contigua (parte está usada por \texttt{ptr1} y
\texttt{ptr2}). La suma de los bloques libres es
$32+128+256+512 = 928$, y aún si fuera mayor, en \emph{buddy system}
\textbf{los pedidos se satisfacen con un único bloque contiguo}.

\begin{importante}
\texttt{ptr3 = malloc(1010)} \textbf{falla} y retorna \texttt{NULL}.
Es un caso clásico de \emph{fragmentación externa}: hay memoria libre,
pero no agrupable en un bloque suficiente.
\end{importante}

El estado del heap no cambia con este pedido fallido.

\subsection{Operación \texttt{free(ptr1)}}

Se libera el bloque de $32$ en el offset $0$. Su \emph{buddy} es el
bloque de $32$ en el offset $32$, que también está libre: se
\textbf{fusionan} en un único bloque de $64$ en el offset $0$.

\textquestiondown{}El nuevo bloque de $64$ tiene un \emph{buddy} libre para seguir
fusionando? Su \emph{buddy} es el bloque de $64$ en el offset $64$
(donde está \texttt{ptr2}), que está \textbf{usado}. La fusión se
detiene.

\begin{center}
\begin{tikzpicture}[
    every node/.style={font=\ttfamily\scriptsize},
    scale=0.85,transform shape,block/.style={draw,minimum height=0.7cm,align=center}
]
\node[block,fill=memfree!20,minimum width=1.0cm] (b0) at (0.5,0) {64};
\node[block,fill=memused!40,minimum width=1.0cm,right=0pt of b0] (b2) {64};
\node[block,fill=memfree!20,minimum width=2.0cm,right=0pt of b2] (b3) {128};
\node[block,fill=memfree!20,minimum width=4.0cm,right=0pt of b3] (b4) {256};
\node[block,fill=memfree!20,minimum width=8.0cm,right=0pt of b4] (b5) {512};
\node[font=\scriptsize] at ($(b0.south)+(0,-0.3)$) {0};
\node[font=\scriptsize] at ($(b5.south east)+(0,-0.3)$) {1024};
\node[font=\bfseries] at ($(b2.north)+(0,0.25)$) {ptr2};
\end{tikzpicture}
\end{center}

\subsection{Operación \texttt{malloc(64)}}

$64 = 2^6$. Hay un bloque libre de exactamente $64$ en el offset $0$,
así que se asigna directamente.

\begin{center}
\begin{tikzpicture}[
    every node/.style={font=\ttfamily\scriptsize},
    scale=0.85,transform shape,block/.style={draw,minimum height=0.7cm,align=center}
]
\node[block,fill=memused!40,minimum width=1.0cm] (b0) at (0.5,0) {64};
\node[block,fill=memused!40,minimum width=1.0cm,right=0pt of b0] (b2) {64};
\node[block,fill=memfree!20,minimum width=2.0cm,right=0pt of b2] (b3) {128};
\node[block,fill=memfree!20,minimum width=4.0cm,right=0pt of b3] (b4) {256};
\node[block,fill=memfree!20,minimum width=8.0cm,right=0pt of b4] (b5) {512};
\node[font=\scriptsize] at ($(b0.south)+(0,-0.3)$) {0};
\node[font=\scriptsize] at ($(b5.south east)+(0,-0.3)$) {1024};
\node[font=\bfseries] at ($(b0.north)+(0,0.25)$) {new};
\node[font=\bfseries] at ($(b2.north)+(0,0.25)$) {ptr2};
\end{tikzpicture}
\end{center}

\subsection*{Resumen}

\begin{center}
\small
\begin{tabular}{@{}lll@{}}
\toprule
Operación & Resultado & Estado del heap (tamaños) \\
\midrule
\texttt{ptr1=malloc(20)}   & OK, asigna 32                & $[\,32_U \mid 32_L \mid 64_L \mid 128_L \mid 256_L \mid 512_L\,]$ \\
\texttt{ptr2=malloc(50)}   & OK, asigna 64                & $[\,32_U \mid 32_L \mid 64_U \mid 128_L \mid 256_L \mid 512_L\,]$ \\
\texttt{ptr3=malloc(1010)} & \textbf{FALLA} (frag.\ ext.) & (sin cambios) \\
\texttt{free(ptr1)}        & libera y fusiona 32+32       & $[\,64_L \mid 64_U \mid 128_L \mid 256_L \mid 512_L\,]$ \\
\texttt{malloc(64)}        & OK, asigna 64                & $[\,64_U \mid 64_U \mid 128_L \mid 256_L \mid 512_L\,]$ \\
\bottomrule
\end{tabular}
\end{center}

\begin{idea}
El \emph{buddy system} simplifica la administración (cada bloque tiene
un único \emph{buddy} fácil de calcular: basta con invertir un bit del
offset) a costa de \emph{fragmentación interna} (cada pedido se
redondea al alza a la siguiente potencia de 2). El ejercicio muestra
ambos costos: \texttt{malloc(20)} consume $32$ bytes (12 desperdiciados,
fragmentación \emph{interna}) y \texttt{malloc(1010)} falla aunque haya
$928$ libres (fragmentación \emph{externa}).
\end{idea}

% =====================================================================
\section{Ejercicio 3 --- Direcciones lógicas de cada área}
% =====================================================================

\textbf{Enunciado.} Hacer un programa en Linux (o MS-Windows) que
permita analizar en qué direcciones lógicas se encuentran las áreas de
código, datos, heap y stack.

\subsection{Programa propuesto}

\begin{lstlisting}[style=cstyle]
#include <stdio.h>
#include <stdlib.h>
#include <unistd.h>

int dato_inicializado = 42;      // .data
int dato_no_inicializado;        // .bss
const char *literal = "hola";    // .rodata (apunta a literal en .text/.rodata)

void funcion(void) { }

int main(int argc, char **argv) {
    int local = 7;                          // stack
    int *heap_ptr = (int *)malloc(sizeof(int));  // heap

    printf("PID                  : %d\n", getpid());
    printf("---- code (.text) ----\n");
    printf("main()               : %p\n", (void *)main);
    printf("funcion()            : %p\n", (void *)funcion);
    printf("literal de cadena    : %p \"%s\"\n", (void *)literal, literal);
    printf("---- data ----\n");
    printf("&dato_inicializado   : %p (.data)\n", (void *)&dato_inicializado);
    printf("&dato_no_inicializado: %p (.bss)\n", (void *)&dato_no_inicializado);
    printf("---- heap ----\n");
    printf("malloc()             : %p\n", (void *)heap_ptr);
    printf("brk actual           : %p\n", sbrk(0));
    printf("---- stack ----\n");
    printf("&local               : %p\n", (void *)&local);
    printf("&argc                : %p\n", (void *)&argc);

    free(heap_ptr);
    return 0;
}
\end{lstlisting}

\subsection{Salida típica}

\begin{lstlisting}[style=textstyle]
PID                  : 18293
---- code (.text) ----
main()               : 0x55a3f4d2b189
funcion()            : 0x55a3f4d2b169
literal de cadena    : 0x55a3f4d2c008 "hola"
---- data ----
&dato_inicializado   : 0x55a3f4d2e010 (.data)
&dato_no_inicializado: 0x55a3f4d2e018 (.bss)
---- heap ----
malloc()             : 0x55a3f5e6c2a0
brk actual           : 0x55a3f5e8d000
---- stack ----
&local               : 0x7fff89a3c2cc
&argc                : 0x7fff89a3c2bc
\end{lstlisting}

\subsection{Lectura del layout}

Las direcciones se ordenan de menor a mayor de la siguiente forma:

\begin{center}
\begin{tikzpicture}[
    every node/.style={font=\ttfamily\small},
    sec/.style={draw,minimum width=4.0cm,minimum height=0.7cm,align=center}
]
\node[sec,fill=mem1!20] (text)  at (0,4.2)  {.text (código)};
\node[sec,fill=mem1!10,below=0pt of text]  (rodata) {.rodata (literales)};
\node[sec,fill=mem2!20,below=0pt of rodata] (data)   {.data (inicializados)};
\node[sec,fill=mem2!10,below=0pt of data]   (bss)    {.bss (no inicializados)};
\node[sec,fill=mem4!20,below=0pt of bss]    (heap)   {heap $\to$ (crece $\uparrow$)};
\node[sec,below=0.8cm of heap]              (gap)    {(espacio libre, librerías, mmap)};
\node[sec,fill=mem3!20,below=0.8cm of gap]  (stack)  {stack ($\downarrow$ crece)};

\node[anchor=west] at ($(text.east)+(0.6,0)$)   {$\sim 0x55\ldots$ (bajo)};
\node[anchor=west] at ($(heap.east)+(0.6,0)$)   {\texttt{brk}};
\node[anchor=west] at ($(stack.east)+(0.6,0)$)  {$\sim 0x7fff\ldots$ (alto)};
\end{tikzpicture}
\end{center}

\begin{itemize}[leftmargin=*,nosep]
    \item \texttt{main}, \texttt{funcion} y los literales (\texttt{.text},
        \texttt{.rodata}) están en la \emph{base baja} del proceso. Son
        de sólo lectura y ejecutables.
    \item \texttt{dato\_inicializado} (\texttt{.data}) y
        \texttt{dato\_no\_inicializado} (\texttt{.bss}) están justo
        después, en la región de datos. \texttt{.bss} ocupa páginas en
        cero rellenadas por el cargador.
    \item \texttt{malloc()} retorna una dirección \emph{por arriba} del
        segmento de datos: es el \textbf{heap}. \texttt{sbrk(0)}
        devuelve el límite superior actual del heap (\emph{program
        break}).
    \item \texttt{\&local} y \texttt{\&argc} están en la dirección
        más alta: el \textbf{stack}, que crece hacia abajo. En Linux
        x86-64, suele empezar cerca de \texttt{0x7fff\ldots}.
\end{itemize}

\subsection{Verificación con \texttt{/proc/<pid>/maps}}

El layout puede contrastarse con la información que el kernel exporta:

\begin{lstlisting}[style=shellstyle]
$ cat /proc/$(pidof a.out)/maps
55a3f4d2b000-55a3f4d2c000 r-xp 00000000 fd:01 ... /home/u/a.out   <- .text
55a3f4d2c000-55a3f4d2d000 r--p 00001000 ...                       <- .rodata
55a3f4d2d000-55a3f4d2e000 r--p 00002000 ...                       <- .data ro
55a3f4d2e000-55a3f4d2f000 rw-p 00003000 ...                       <- .data rw + .bss
55a3f5e6b000-55a3f5e8d000 rw-p 00000000 00:00 0   [heap]
7f...     -7f...           r-xp ...   libc.so.6
7ffe89a1c000-7ffe89a3d000 rw-p 00000000 00:00 0   [stack]
\end{lstlisting}

\begin{idea}
La maquinaria de paginado le permite al sistema dejar \emph{huecos}
gigantes entre \texttt{heap} y \texttt{stack}, sin reservar memoria
física por ellos: los \emph{pte} intermedios son inválidos
($V=0$) y sólo se materializan en cuanto el proceso los referencie.
Esto es lo que viabiliza el crecimiento dinámico del stack y del heap.
\end{idea}

% =====================================================================
\section{Ejercicio 4 --- Arquitectura de 16 bits con páginas de 1 KB}
% =====================================================================

\textbf{Enunciado.} Arquitectura con espacio de direcciones de
$16$\,bits, paginado con páginas de $1$\,KB.

\subsection{4.a --- Capacidad máxima de la memoria}

Con direcciones de $16$\,bits, el espacio direccionable tiene
$2^{16} = 65\,536$ bytes, es decir:
\[
\boxed{\textit{capacidad máxima} = 2^{16}\,\text{bytes} = 64\,\text{KB}.}
\]

\subsection{4.b --- Cantidad de páginas}

Con páginas de $1\,\text{KB} = 2^{10}$ bytes:
\[
\textit{número de páginas} = \frac{2^{16}}{2^{10}} = 2^{6} = \boxed{64\ \text{páginas}.}
\]

\subsection{4.c --- Representación de la tabla de páginas y de la \emph{pte}}

Una dirección lógica de $16$\,bits se divide en:

\begin{center}
\begin{tikzpicture}[
    every node/.style={font=\ttfamily\small},
    cell/.style={draw,minimum height=0.7cm,align=center}
]
\node[cell,minimum width=2.2cm,fill=mem1!20] (vpn) at (0,0) {vpn};
\node[cell,minimum width=4.0cm,fill=mem2!20,right=0pt of vpn] (off) {offset};
\node[font=\scriptsize] at ($(vpn.north)+(0,0.25)$) {6 bits};
\node[font=\scriptsize] at ($(off.north)+(0,0.25)$) {10 bits};
\node[font=\scriptsize] at ($(vpn.north west)+(0.05,-0.05)$) {15};
\node[font=\scriptsize] at ($(vpn.north east)+(-0.1,-0.05)$) {10};
\node[font=\scriptsize] at ($(off.north west)+(0.05,-0.05)$) {9};
\node[font=\scriptsize] at ($(off.north east)+(-0.05,-0.05)$) {0};
\end{tikzpicture}
\end{center}

\begin{itemize}[leftmargin=*,nosep]
    \item \texttt{vpn} ($6$ bits): índice en la tabla de páginas
        ($2^{6} = 64$ entradas).
    \item \texttt{offset} ($10$ bits): desplazamiento dentro de la
        página de $1\,\text{KB}$.
\end{itemize}

\paragraph{Formato propuesto de \emph{pte}.} Cada entrada describe
dónde está la página física y con qué permisos. Un diseño compacto de
$16$\,bits sería:

\begin{center}
\begin{tikzpicture}[
    every node/.style={font=\ttfamily\small},
    cell/.style={draw,minimum height=0.7cm,align=center}
]
\node[cell,minimum width=2.4cm,fill=mem1!20] (ppn) at (0,0) {ppn};
\node[cell,minimum width=0.6cm,fill=mem3!30,right=0pt of ppn] (V) {V};
\node[cell,minimum width=0.6cm,fill=mem3!30,right=0pt of V] (R) {R};
\node[cell,minimum width=0.6cm,fill=mem3!30,right=0pt of R] (W) {W};
\node[cell,minimum width=0.6cm,fill=mem3!30,right=0pt of W] (X) {X};
\node[cell,minimum width=0.6cm,fill=mem3!30,right=0pt of X] (U) {U};
\node[cell,minimum width=0.6cm,fill=mem3!30,right=0pt of U] (A) {A};
\node[cell,minimum width=0.6cm,fill=mem3!30,right=0pt of A] (D) {D};
\node[cell,minimum width=1.0cm,right=0pt of D] (res) {res.};
\node[font=\scriptsize] at ($(ppn.north)+(0,0.25)$) {6 bits};
\node[font=\scriptsize] at ($(V.north)+(0,0.25)$) {1};
\node[font=\scriptsize] at ($(R.north)+(0,0.25)$) {1};
\node[font=\scriptsize] at ($(W.north)+(0,0.25)$) {1};
\node[font=\scriptsize] at ($(X.north)+(0,0.25)$) {1};
\node[font=\scriptsize] at ($(U.north)+(0,0.25)$) {1};
\node[font=\scriptsize] at ($(A.north)+(0,0.25)$) {1};
\node[font=\scriptsize] at ($(D.north)+(0,0.25)$) {1};
\node[font=\scriptsize] at ($(res.north)+(0,0.25)$) {3};
\end{tikzpicture}
\end{center}

Los campos son los habituales (cf.\ resumen teórico de gestión de
memoria, formato \emph{pte} de RV-32):

\begin{itemize}[leftmargin=*,nosep]
    \item \texttt{ppn} ($6$\,bits): número de página física
        (\emph{physical page number}).
    \item \texttt{V}: bit \emph{valid}. Si $V=0$, indexar esa entrada
        provoca un \emph{page fault}.
    \item \texttt{R}, \texttt{W}, \texttt{X}: permisos de \emph{read},
        \emph{write} y \emph{execute}.
    \item \texttt{U}: permiso en modo \emph{user}.
    \item \texttt{A}: encendido por el hardware cuando se accede
        (lectura o escritura).
    \item \texttt{D}: encendido por el hardware cuando se escribe
        (\emph{dirty}).
\end{itemize}

\paragraph{Tamaño total de la tabla.} $64$ entradas $\times$ $16$\,bits
$= 1024$\,bits $= 128$\,bytes, lo que cabe holgadamente en una sola
página de $1$\,KB.

\subsection{4.d --- Tabla de páginas de un proceso concreto}

El proceso tiene:

\begin{itemize}[leftmargin=*,nosep]
    \item $1800$\,bytes de \textbf{código};
    \item $600$\,bytes de \textbf{datos} (a continuación);
    \item una \textbf{página de guarda} entre datos y pila (para
        detectar \emph{stack overflow});
    \item $1$\,KB de \textbf{pila}.
\end{itemize}

\paragraph{Cuántas páginas usa cada área.}

\begin{align*}
\text{código} &: \lceil 1800/1024 \rceil = 2\ \text{páginas} & \text{(páginas 0 y 1)} \\
\text{datos}  &: \lceil 600/1024 \rceil  = 1\ \text{página} & \text{(página 2)} \\
\text{guarda} &: \text{1 página inválida}                     & \text{(página 3)} \\
\text{pila}   &: \lceil 1024/1024 \rceil = 1\ \text{página} & \text{(página 4)}
\end{align*}

El espacio de direcciones del proceso queda $[0, 5\,\text{KB}]$. Las
entradas $5$ a $63$ no se usan (\texttt{V}=0).

\paragraph{Tabla de páginas.} Los \texttt{ppn} concretos los elige el
SO según los \emph{frames} libres; en lo que sigue se usan
\texttt{0x07}, \texttt{0x12}, \texttt{0x1A} y \texttt{0x2F} a modo de
ejemplo.

\begin{center}
\small
\begin{tabular}{@{}clccccccll@{}}
\toprule
\#  & Área           & \texttt{ppn}  & V & R & W & X & U & D/A & Comentario \\
\midrule
0   & código        & \texttt{0x07} & 1 & 1 & 0 & 1 & 1 & --  & primera página de \texttt{.text} \\
1   & código        & \texttt{0x12} & 1 & 1 & 0 & 1 & 1 & --  & segunda página de \texttt{.text} \\
2   & datos         & \texttt{0x1A} & 1 & 1 & 1 & 0 & 1 & --  & \texttt{.data} + \texttt{.bss} \\
3   & \textbf{guarda} & ---         & 0 & 0 & 0 & 0 & 0 & --  & cualquier acceso $\Rightarrow$ \emph{page fault} \\
4   & pila          & \texttt{0x2F} & 1 & 1 & 1 & 0 & 1 & --  & 1\,KB de stack \\
5\,--\,63 & ---     & ---           & 0 & 0 & 0 & 0 & 0 & --  & páginas no asignadas \\
\bottomrule
\end{tabular}
\end{center}

\begin{center}
\begin{tikzpicture}[
    every node/.style={font=\ttfamily\small},
    sec/.style={draw,minimum width=3.5cm,minimum height=0.7cm,align=center}
]
\node[sec,fill=mem1!20]  (c0) at (0,5.6) {página 0 --- código};
\node[sec,fill=mem1!20,below=0pt of c0] (c1) {página 1 --- código};
\node[sec,fill=mem2!20,below=0pt of c1] (d)  {página 2 --- datos};
\node[sec,fill=memused!30,below=0pt of d] (g) {página 3 --- guarda (V=0)};
\node[sec,fill=mem3!20,below=0pt of g] (s)  {página 4 --- pila};
\node[anchor=west,xshift=4pt] at (c0.east) {\texttt{0x0000}};
\node[anchor=west,xshift=4pt] at (c1.east) {\texttt{0x0400}};
\node[anchor=west,xshift=4pt] at (d.east)  {\texttt{0x0800}};
\node[anchor=west,xshift=4pt] at (g.east)  {\texttt{0x0C00}};
\node[anchor=west,xshift=4pt] at (s.east)  {\texttt{0x1000}};
\node[anchor=west,xshift=4pt] at ($(s.south east)+(0,0)$) {\texttt{0x1400}};
\end{tikzpicture}
\end{center}

\begin{idea}
Si el stack crece más allá de la página 4, la siguiente referencia cae
en la página 3 (guarda, $V=0$) y el hardware genera un \emph{page
fault}. El kernel reconoce el patrón ``escritura a una página de
guarda inválida'' y aborta el proceso con \emph{stack overflow}. Es
exactamente el mecanismo descripto en el resumen teórico de gestión
de memoria.
\end{idea}

% =====================================================================
\section{Ejercicio 5 --- Tablas de páginas para 4 GB en 32 bits}
% =====================================================================

\textbf{Enunciado.} En un SO sobre una arquitectura de $32$\,bits con
paginado, determinar la cantidad de \emph{ptes} y la memoria total
necesaria para cubrir $4$\,GB de espacio de direcciones, con
(a) páginas de $4$\,KB y (b) páginas de $4$\,MB.

Se asume el tamaño habitual de cada \emph{pte}: $32$\,bits
$=4$\,bytes (espacio para el \texttt{ppn} más los flags). El cálculo
general es
\[
\#\text{ptes} = \frac{2^{32}}{\text{tamaño página}}, \qquad
\text{memoria de la tabla} = \#\text{ptes} \times 4\,\text{bytes}.
\]

\subsection{5.a --- Páginas de 4 KB}

\begin{align*}
\#\text{ptes} &= \frac{2^{32}}{2^{12}} = 2^{20} = 1\,048\,576 \\
\text{memoria de la tabla} &= 2^{20} \times 4\,\text{B} = 2^{22}\,\text{B} = 4\,\text{MB}.
\end{align*}

Una tabla \emph{plana} de $4$\,MB \textbf{por proceso} es enorme y
exigiría un bloque contiguo de RAM grande para alojarla. Por eso los
SO modernos representan la tabla como un \textbf{árbol de varios
niveles} (por ejemplo, 2 niveles en RV-32, 4 niveles en x86-64): sólo
se reservan los nodos correspondientes a regiones del espacio de
direcciones efectivamente usadas, y los enormes huecos no consumen
memoria.

\subsection{5.b --- Páginas de 4 MB}

\begin{align*}
\#\text{ptes} &= \frac{2^{32}}{2^{22}} = 2^{10} = 1024 \\
\text{memoria de la tabla} &= 2^{10} \times 4\,\text{B} = 2^{12}\,\text{B} = 4\,\text{KB}.
\end{align*}

Una tabla de sólo $4$\,KB cabe en una sola página, lo que simplifica
muchísimo la administración. El costo es la \textbf{fragmentación
interna}: cualquier asignación de memoria, por pequeña que sea, ocupa
$4$\,MB completos.

\subsection*{Resumen comparativo}

\begin{center}
\small
\begin{tabular}{@{}lcccc@{}}
\toprule
Caso & Tamaño de página & \#ptes & Memoria de la tabla & Comentario \\
\midrule
(a) & $4$\,KB   & $2^{20} = 1.048.576$ & $4$\,MB & menor frag.\ interna; tabla enorme \\
(b) & $4$\,MB   & $2^{10} = 1.024$     & $4$\,KB & tabla chica; gran frag.\ interna \\
\bottomrule
\end{tabular}
\end{center}

\begin{importante}
Hay un compromiso fundamental: \textbf{páginas más chicas} dan menor
fragmentación interna pero tablas mucho más grandes; \textbf{páginas
más grandes} (\emph{huge pages}) reducen el tamaño de la tabla y la
presión sobre el TLB, pero malgastan memoria si los pedidos son
pequeños. Los SO modernos combinan ambos: páginas de $4$\,KB por
default, con la opción de usar páginas de $2$\,MB / $1$\,GB
(\emph{huge pages}) en regiones específicas.
\end{importante}

% =====================================================================
\section{Ejercicio 6 --- Representación de la tabla de páginas de un proceso}
% =====================================================================

\textbf{Enunciado.} Representar la tabla de páginas de un proceso.

Se toma el proceso del Ejercicio 4d (el ejercicio del práctico no fija
un proceso particular, así que se reutiliza ese ejemplo para
concretar). La representación enfatiza la traducción de una dirección
lógica a una dirección física a través de la \emph{pte} indexada.

\begin{center}
\begin{tikzpicture}[
    scale=0.85,transform shape,
    every node/.style={font=\ttfamily\small},
    cell/.style={draw,minimum height=0.65cm,align=center},
    page/.style={draw,minimum width=1.6cm,minimum height=0.65cm,align=center}
]

% Logical addr
\node[font=\bfseries\small] at (-1.3,4.8) {dirección lógica};
\node[cell,minimum width=1.0cm,fill=mem1!20] (vpn) at (-2.1,3.6) {vpn};
\node[cell,minimum width=2.0cm,fill=mem2!20,right=0pt of vpn] (off) {offset};
\node[font=\scriptsize] at ($(vpn.north)+(0,0.25)$) {6 b};
\node[font=\scriptsize] at ($(off.north)+(0,0.25)$) {10 b};

% Page table
\node[font=\bfseries\small] at (3.5,4.8) {tabla de páginas};
\node[cell,minimum width=1.2cm,fill=mem1!10] (e0p) at (3.0,3.6) {\texttt{0x07}};
\node[cell,minimum width=2.0cm,fill=mem3!20,right=0pt of e0p] (e0f) {V=1 R=1 X=1};
\node[anchor=east] at (e0p.west) {0};

\node[cell,minimum width=1.2cm,fill=mem1!10] (e1p) at (3.0,2.9) {\texttt{0x12}};
\node[cell,minimum width=2.0cm,fill=mem3!20,right=0pt of e1p] (e1f) {V=1 R=1 X=1};
\node[anchor=east] at (e1p.west) {1};

\node[cell,minimum width=1.2cm,fill=mem1!10] (e2p) at (3.0,2.2) {\texttt{0x1A}};
\node[cell,minimum width=2.0cm,fill=mem3!20,right=0pt of e2p] (e2f) {V=1 R=1 W=1};
\node[anchor=east] at (e2p.west) {2};

\node[cell,minimum width=1.2cm,fill=memused!20] (e3p) at (3.0,1.5) {---};
\node[cell,minimum width=2.0cm,fill=memused!20,right=0pt of e3p] (e3f) {V=0 (guarda)};
\node[anchor=east] at (e3p.west) {3};

\node[cell,minimum width=1.2cm,fill=mem1!10] (e4p) at (3.0,0.8) {\texttt{0x2F}};
\node[cell,minimum width=2.0cm,fill=mem3!20,right=0pt of e4p] (e4f) {V=1 R=1 W=1};
\node[anchor=east] at (e4p.west) {4};

\node[cell,minimum width=3.2cm] (edots) at (3.6,0.1) {\ldots};
\node[anchor=east] at (edots.west) {5--63};

% Physical memory
\node[font=\bfseries\small] at (9.5,4.3) {memoria física};
\node[page,fill=mem1!20] (f07) at (9.5,3.6) {frame \texttt{0x07}};
\node[page,fill=mem1!20] (f12) at (9.5,2.9) {frame \texttt{0x12}};
\node[page,fill=mem2!20] (f1A) at (9.5,2.2) {frame \texttt{0x1A}};
\node[page,fill=mem3!20] (f2F) at (9.5,0.8) {frame \texttt{0x2F}};

% Arrows
\draw[-Latex,thick] (vpn.south) -- ++(0,-0.4) -| (e0p.west);
\draw[-Latex,thick] (e0f.east) -- (f07.west);
\draw[-Latex,thick] (e1f.east) -- (f12.west);
\draw[-Latex,thick] (e2f.east) -- (f1A.west);
\draw[-Latex,thick] (e4f.east) -- (f2F.west);
\end{tikzpicture}
\end{center}

\paragraph{Traducción de una dirección lógica.}
Dada la dirección lógica $\textit{va}$, se separa en
$(\textit{vpn}, \textit{offset})$ con los $6$ bits altos y $10$ bits
bajos. Se accede a la entrada $\textit{vpn}$ de la tabla. Si $V=1$ y
los permisos cuadran con el tipo de acceso, la dirección física
queda
\[
\textit{pa} = \texttt{pagetable}[\textit{vpn}].\texttt{ppn}
              \times 1024 + \textit{offset}.
\]
Si $V=0$ o se viola algún permiso, el hardware lanza un \emph{page
fault}.

\paragraph{Ejemplo numérico.} La instrucción referencia la dirección
lógica \texttt{0x0814} (en el área de datos). Se descompone como
$\textit{vpn} = 2$ y $\textit{offset} = \texttt{0x014} = 20$. La
\emph{pte} 2 tiene $\texttt{ppn} = \texttt{0x1A}$ y $V=1$, $R=1$,
$W=1$. La dirección física es
\[
\texttt{0x1A} \times 1024 + 20 = 26\,624 + 20 = 26\,644 = \texttt{0x6814}.
\]

% =====================================================================
\section{Ejercicio 7 --- Memoria compartida entre dos procesos}
% =====================================================================

\textbf{Enunciado.} Dos procesos comparten un área de memoria de $2$
páginas. En el primer proceso el área se mapea a partir de
\texttt{0xA000}; en el segundo desde \texttt{0xF000}. Mostrar con un
diagrama el layout de memoria de cada proceso y las tablas de páginas.
Las páginas son de $4$\,KB.

\subsection{Análisis de las direcciones lógicas}

Con páginas de $4\,\text{KB} = 2^{12}$, el \emph{offset} ocupa los
$12$ bits menos significativos y el \texttt{vpn} ocupa los bits altos.
Por lo tanto:

\begin{center}
\small
\begin{tabular}{@{}lll@{}}
\toprule
Dirección & vpn & offset \\
\midrule
\texttt{0xA000} & $\texttt{0xA} = 10$ & $0$ \\
\texttt{0xB000} & $\texttt{0xB} = 11$ & $0$ \\
\texttt{0xF000} & $\texttt{0xF} = 15$ & $0$ \\
\texttt{0x10000} & $\texttt{0x10} = 16$ & $0$ \\
\bottomrule
\end{tabular}
\end{center}

\textbf{Proceso 1.} El área compartida ocupa las páginas lógicas
$10$ y $11$ (rango $[\texttt{0xA000}, \texttt{0xC000})$).
\textbf{Proceso 2.} El área compartida ocupa las páginas lógicas
$15$ y $16$ (rango $[\texttt{0xF000}, \texttt{0x11000})$).

En ambos procesos, esas \emph{ptes} apuntan a los \textbf{mismos dos
frames físicos} $X$ e $Y$.

\subsection{Diagrama}

\begin{center}
\begin{tikzpicture}[
    every node/.style={font=\ttfamily\small},
    sec/.style={draw,minimum width=3.6cm,minimum height=0.6cm,align=center},
    fr/.style={draw,minimum width=2.0cm,minimum height=0.6cm,align=center}
]

% Proceso 1 layout
\node[font=\bfseries\small] at (-4,5.5) {proceso 1};
\node[sec,fill=mem1!20] (p1_text) at (-4,4.8) {.text};
\node[sec,fill=mem2!20,below=0pt of p1_text] (p1_data) {.data};
\node[sec,below=0pt of p1_data] (p1_gap1) {\ldots};
\node[sec,fill=memshared!50,below=0pt of p1_gap1] (p1_sh1) {compartida \#0};
\node[sec,fill=memshared!50,below=0pt of p1_sh1] (p1_sh2) {compartida \#1};
\node[sec,below=0pt of p1_sh2] (p1_gap2) {\ldots};
\node[sec,fill=mem3!20,below=0pt of p1_gap2] (p1_stack) {stack};
\node[anchor=east] at (p1_sh1.west) {\texttt{0xA000}};
\node[anchor=east] at (p1_sh2.west) {\texttt{0xB000}};
\node[anchor=east] at ($(p1_sh2.south west)+(0,-0.05)$) {\texttt{0xC000}};

% Proceso 2 layout
\node[font=\bfseries\small] at (4,5.5) {proceso 2};
\node[sec,fill=mem1!20] (p2_text) at (4,4.8) {.text};
\node[sec,fill=mem2!20,below=0pt of p2_text] (p2_data) {.data};
\node[sec,below=0pt of p2_data] (p2_gap1) {\ldots};
\node[sec,fill=memshared!50,below=0pt of p2_gap1] (p2_sh1) {compartida \#0};
\node[sec,fill=memshared!50,below=0pt of p2_sh1] (p2_sh2) {compartida \#1};
\node[sec,below=0pt of p2_sh2] (p2_gap2) {\ldots};
\node[sec,fill=mem3!20,below=0pt of p2_gap2] (p2_stack) {stack};
\node[anchor=west] at (p2_sh1.east) {\texttt{0xF000}};
\node[anchor=west] at (p2_sh2.east) {\texttt{0x10000}};
\node[anchor=west] at ($(p2_sh2.south east)+(0,-0.05)$) {\texttt{0x11000}};

% Shared frames in the middle
\node[font=\bfseries\small] at (0,2.5) {memoria física};
\node[fr,fill=memshared!60] (fX) at (0,1.8) {frame X};
\node[fr,fill=memshared!60,below=0pt of fX] (fY) {frame Y};

% Arrows from each process to the shared frames
\draw[-Latex,thick] (p1_sh1.east) -- ++(0.8,0) |- (fX.west);
\draw[-Latex,thick] (p1_sh2.east) -- ++(0.4,0) |- (fY.west);
\draw[-Latex,thick] (p2_sh1.west) -- ++(-0.8,0) |- (fX.east);
\draw[-Latex,thick] (p2_sh2.west) -- ++(-0.4,0) |- (fY.east);

\end{tikzpicture}
\end{center}

\subsection{Tablas de páginas}

\begin{center}
\small
\textbf{Proceso 1} \\[0.3em]
\begin{tabular}{@{}clcccc@{}}
\toprule
vpn & decimal & ppn & V & R & W \\
\midrule
\ldots & \ldots & \ldots & \ldots & \ldots & \ldots \\
\texttt{0xA} & 10 & $X$ & 1 & 1 & 1 \\
\texttt{0xB} & 11 & $Y$ & 1 & 1 & 1 \\
\ldots & \ldots & \ldots & \ldots & \ldots & \ldots \\
\bottomrule
\end{tabular}
\qquad
\textbf{Proceso 2} \\[0.3em]
\begin{tabular}{@{}clcccc@{}}
\toprule
vpn & decimal & ppn & V & R & W \\
\midrule
\ldots & \ldots & \ldots & \ldots & \ldots & \ldots \\
\texttt{0xF}  & 15 & $X$ & 1 & 1 & 1 \\
\texttt{0x10} & 16 & $Y$ & 1 & 1 & 1 \\
\ldots & \ldots & \ldots & \ldots & \ldots & \ldots \\
\bottomrule
\end{tabular}
\end{center}

\begin{idea}
La clave del esquema es que \emph{distintas} entradas
($\texttt{vpn}=10$ en P1, $\texttt{vpn}=15$ en P2) almacenan
\emph{el mismo} \texttt{ppn}. Las direcciones lógicas son distintas
en cada proceso, pero el hardware traduce ambas al mismo \emph{frame}
físico. Así, una escritura del proceso 1 a \texttt{0xA000} es
inmediatamente visible para el proceso 2 leyendo \texttt{0xF000}. Es
exactamente el mecanismo que usan \texttt{shmat} y \texttt{mmap} con
\texttt{MAP\_SHARED}.
\end{idea}

\begin{importante}
La memoria compartida no es un canal sincronizado por sí mismo: si
ambos procesos escriben simultáneamente al área compartida, pueden
producirse \emph{race conditions}. Es responsabilidad del programador
agregar un mecanismo de sincronización (semáforos, mutex en memoria
compartida, \texttt{futex}, etc.).
\end{importante}

% =====================================================================
\section{Ejercicio 8 --- mmap de un archivo}
% =====================================================================

\textbf{Enunciado.} Compilar y ejecutar (en Linux) el archivo
\texttt{mmapfile.c} dado. (a) \textquestiondown{}Por qué \texttt{data + 7000} no causa
\emph{segmentation fault}, si el archivo ocupa poco más de $4096$
bytes? (b) Modificar el programa para mapear en lectura y escritura,
con persistencia automática al \texttt{munmap}; \textquestiondown{}cómo lo hace el SO?
(c) Verificar el tamaño del archivo después de ejecutar.

\subsection{Listado del programa original}

\begin{lstlisting}[style=cstyle]
#include <stdio.h>
#include <string.h>
#include <fcntl.h>
#include <unistd.h>
#include <sys/mman.h>

#define PGSIZE   4096
#define filename "mmapfile.txt"

int createfile(void) {
    int f = open(filename, O_RDONLY);
    if (f < 0) {
        const char *str1 = "Hello world in page 1";
        const char *str2 = "I'm in second page";
        const unsigned char zero = 0;
        int l = strlen(str1), i;
        f = open(filename, O_RDWR | O_CREAT, 0666);
        write(f, str1, l);
        for (i = l; i < PGSIZE; i++) write(f, &zero, 1);
        write(f, str2, strlen(str2) + 1);
        close(f);
        f = open(filename, O_RDONLY);
    }
    return f;
}

int main(void) {
    int f = createfile();
    char *data = mmap(0, PGSIZE * 2, PROT_READ, MAP_PRIVATE, f, 0);
    if (data) {
        printf("file contents mapped at %p\n", data);
        printf("%s\n", data);
        printf("%s\n", data + 4096);
        printf("%s\n", data + 7000); /* acceso valido aunque el archivo "termine" antes */
        munmap(data, PGSIZE * 2);
    } else {
        printf("mmap error\n");
    }
    close(f);
    return 0;
}
\end{lstlisting}

\subsection{Compilación y ejecución}

\begin{lstlisting}[style=shellstyle]
$ gcc -Wall -O0 -o mmapfile mmapfile.c
$ ./mmapfile
file contents mapped at 0x7f8e2c1c4000
Hello world in page 1
I'm in second page
                       <- aparece linea en blanco (los bytes son ceros)
\end{lstlisting}

\subsection{8.a --- \textquestiondown{}Por qué \texttt{data + 7000} no causa segfault?}

El archivo \texttt{mmapfile.txt} pesa $\texttt{strlen(str1)} + (PGSIZE
- \texttt{strlen(str1)}) + \texttt{strlen(str2)} + 1 \approx 4115$
bytes, es decir, poco más de una página.

Sin embargo, la llamada \verb|mmap(0, PGSIZE * 2, ..., f, 0)| pidió
mapear \textbf{dos} páginas ($8192$ bytes) a partir del offset $0$.

\begin{importante}
La región mapeada es \emph{redondeada hacia arriba al tamaño de
página}. Por eso \texttt{mmap} reserva 2 páginas válidas en el espacio
de direcciones, aunque el archivo no llene la segunda completamente.
\end{importante}

El kernel, en ese caso, \textbf{rellena con ceros} la parte de la
última página que excede al EOF del archivo, según lo definido por
POSIX. Por lo tanto, todas las direcciones del rango
$[\texttt{data}, \texttt{data} + 8192)$ son válidas y, en particular,
$\texttt{data} + 7000$:

\begin{itemize}[leftmargin=*,nosep]
    \item Está \emph{dentro} de la región mapeada (offsets
        $[0, 8192)$): la \texttt{pte} correspondiente está definida.
    \item Los bytes leídos a partir de allí son los del archivo en el
        offset $7000 - 4096 = 2904$ de la página $1$. Como el archivo
        no llega hasta ahí, son \textbf{ceros} (rellenados por el
        kernel).
\end{itemize}

Por eso \texttt{printf("\%s\textbackslash{}n", data + 7000)} imprime
una cadena vacía (el primer byte es $\backslash 0$) en lugar de
provocar un \emph{segmentation fault}.

\begin{idea}
Si en cambio se accediese a \texttt{data + 8192} (la página $2$, no
mapeada), la \texttt{pte} no existiría y sí se produciría
\emph{segfault} (o más precisamente, \emph{SIGSEGV}). El mecanismo
exacto es el visto en el resumen teórico de memoria virtual: el
hardware levanta un \emph{page fault}, el kernel detecta que la
dirección no pertenece a ningún VMA del proceso y, al no ser
recuperable, envía la señal correspondiente.
\end{idea}

\subsection{8.b --- Modificación para lectura y escritura con persistencia}

Hay que cambiar dos cosas en la llamada a \texttt{mmap}:

\begin{enumerate}[leftmargin=*,nosep]
    \item Permisos de la página: \verb!PROT_READ! $\to$ \verb!PROT_READ | PROT_WRITE!.
    \item Modo de compartición: \verb!MAP_PRIVATE! $\to$ \verb!MAP_SHARED!.
\end{enumerate}

\begin{importante}
Con \texttt{MAP\_PRIVATE} las escrituras al área mapeada generan una
copia privada (\emph{copy on write}) y \emph{no} se propagan al
archivo. Para que las modificaciones queden en el archivo se necesita
\texttt{MAP\_SHARED}. Además, el archivo debe haberse abierto en
modo \texttt{O\_RDWR}.
\end{importante}

\begin{lstlisting}[style=cstyle]
int main(void) {
    int f = createfile_rw();    // abre con O_RDWR | O_CREAT
    char *data = mmap(0, PGSIZE * 2, PROT_READ | PROT_WRITE,
                      MAP_SHARED, f, 0);
    if (data == MAP_FAILED) { perror("mmap"); return 1; }

    printf("file contents mapped at %p\n", data);
    printf("Before: %s\n", data);

    /* modifico el contenido a traves de la memoria mapeada */
    const char *nuevo = "Hola desde mmap()!";
    memcpy(data, nuevo, strlen(nuevo) + 1);

    /* munmap se encarga de hacer flush de las paginas dirty al archivo */
    munmap(data, PGSIZE * 2);
    close(f);
    return 0;
}
\end{lstlisting}

\paragraph{Cómo lo hace el SO.}
La explicación está en el resumen teórico de memoria virtual:

\begin{enumerate}[leftmargin=*]
    \item Por cada página mapeada, el hardware mantiene el bit
        \textbf{\texttt{D}} (\emph{dirty}) en la \texttt{pte}. Cada
        vez que el proceso escribe en esa página, el \texttt{D} pasa
        a $1$ automáticamente.
    \item Cuando se invoca \texttt{munmap} (o cuando el sistema hace
        un \emph{flush} periódico, p.\,ej.\ \texttt{pdflush}), el
        kernel recorre las \texttt{pte} del mapping y escribe al
        archivo \emph{sólo las páginas con \texttt{D}=1}, en el
        offset correspondiente. Las páginas con \texttt{D}=0 no se
        tocan, ahorrando E/S.
    \item Finalmente, la región del espacio de direcciones se libera
        y las \texttt{pte} pasan a $V=0$.
\end{enumerate}

Adicionalmente, la llamada \texttt{msync(addr, len, MS\_SYNC)} permite
forzar la sincronización a disco \emph{sin} desmontar el mapping; es
útil cuando se quiere garantizar persistencia antes del \texttt{munmap}
(por ejemplo, frente a un eventual crash).

\subsection{8.c --- Tamaño del archivo después}

\begin{lstlisting}[style=shellstyle]
$ ls -l mmapfile.txt
-rw-r--r-- 1 tomi tomi 4115 may 20 23:01 mmapfile.txt
$ stat -c "size: %s bytes" mmapfile.txt
size: 4115 bytes
\end{lstlisting}

\textbf{El tamaño no cambia.} Aunque el mapping cubre dos páginas
($8192$ bytes), el archivo subyacente sigue ocupando los $\sim 4115$
bytes que tenía antes. Las escrituras realizadas en offsets
$\geq 4115$ \emph{dentro de la primera página} se propagan a esa
zona del archivo siempre que no excedan el tamaño actual; las
escrituras dentro de la segunda página \emph{sí} podrían propagarse,
pero el comportamiento ante escrituras más allá del EOF en un mapping
de un archivo regular no está definido por POSIX (en Linux suele
generar \texttt{SIGBUS}).

\begin{importante}
Para extender el archivo, hay que llamar a \texttt{ftruncate(f,
new\_size)} \emph{antes} de \texttt{mmap}. Sólo entonces la región
mapeada coincide con un archivo de ese tamaño y las escrituras quedan
persistidas correctamente.
\end{importante}

% =====================================================================
\section{Ejercicio 9 --- Memoria compartida IPC con shmwriter y shmreader}
% =====================================================================

\textbf{Enunciado.} Compilar y ejecutar (en ese orden)
\texttt{shmwriter.c} y \texttt{shmreader.c} para observar cómo
funciona la API de \emph{shared memory} (IPC de UNIX). (a) Volver a
correr \texttt{shmreader} y ver qué ocurre. (b) Comentar la línea
\texttt{shmctl(..., IPC\_RMID, NULL)} de \texttt{shmreader.c}; correr
\texttt{shmwriter} y luego \texttt{shmreader} varias veces.
(c) Modificar ambos programas para leer/escribir en la dirección
lógica \texttt{str+500}.

\subsection{Cómo funciona la API}

\begin{itemize}[leftmargin=*]
    \item \verb|key_t ftok(path, id)|: deriva una clave numérica a
        partir de un archivo existente y un \texttt{id} (un byte). Los
        procesos cooperantes acuerdan el mismo \texttt{path} y
        \texttt{id} para coincidir en la misma clave.
    \item \verb|int shmget(key, size, flags)|: obtiene (o crea con
        \texttt{IPC\_CREAT}) un segmento de \emph{shared memory} con
        esa clave y tamaño. Devuelve un identificador \texttt{shmid}.
    \item \verb|void *shmat(shmid, addr, flags)|: mapea el segmento
        en el espacio de direcciones del proceso. Si \texttt{addr=0},
        el kernel elige la dirección.
    \item \verb|int shmdt(addr)|: hace \emph{detach}: quita el mapping
        de este proceso, pero el segmento sigue existiendo.
    \item \verb|int shmctl(shmid, IPC_RMID, NULL)|: marca el segmento
        para destrucción. Se libera cuando el último proceso hace
        \texttt{shmdt}.
\end{itemize}

\subsection{Programas provistos}

\begin{lstlisting}[style=cstyle,caption={shmwriter.c}]
#include <sys/ipc.h>
#include <sys/shm.h>
#include <stdio.h>
#include <string.h>

int main(void) {
    key_t key = ftok("shmfile", 65);
    int shmid = shmget(key, 1024, 0666 | IPC_CREAT);
    char *str = (char *)shmat(shmid, (void *)0, 0);

    printf("shared memory logical address: %p\n", str);
    printf("Data to write to shmem: ");
    fgets(str, 1024, stdin);
    str[strcspn(str, "\n")] = '\0';
    printf("Data written in memory: %s\n", str);

    shmdt(str);
    return 0;
}
\end{lstlisting}

\begin{lstlisting}[style=cstyle,caption={shmreader.c}]
#include <sys/ipc.h>
#include <sys/shm.h>
#include <stdio.h>

int main(void) {
    key_t key = ftok("shmfile", 65);
    int shmid = shmget(key, 1024, 0666 | IPC_CREAT);
    char *str = (char *)shmat(shmid, (void *)0, 0);

    printf("shared memory logical address: %p\n", str);
    printf("Data read from memory: %s\n", str);

    shmdt(str);
    shmctl(shmid, IPC_RMID, NULL);  /* destruye el segmento */
    return 0;
}
\end{lstlisting}

\subsection{Ejecución base}

\begin{lstlisting}[style=shellstyle]
$ touch shmfile
$ gcc -o shmwriter shmwriter.c
$ gcc -o shmreader shmreader.c

$ ./shmwriter
shared memory logical address: 0x7fbcde47b000
Data to write to shmem: Hola desde shmwriter
Data written in memory: Hola desde shmwriter

$ ./shmreader
shared memory logical address: 0x7f0123abc000
Data read from memory: Hola desde shmwriter
\end{lstlisting}

\textbf{Observación.} Las direcciones lógicas reportadas son
\emph{diferentes} ($\texttt{0x7fbc\ldots}$ vs $\texttt{0x7f01\ldots}$),
pero los datos coinciden: ambos procesos están viendo los mismos
\emph{frames} físicos a través de \texttt{vpn} distintos. Es
exactamente el escenario del Ejercicio 7.

\subsection{9.a --- Volver a correr \texttt{shmreader}}

Si \texttt{shmreader} \emph{destruye} el segmento al terminar (con
\texttt{shmctl(IPC\_RMID)}, como está escrito arriba), la segunda
ejecución encuentra el siguiente escenario:

\begin{itemize}[leftmargin=*,nosep]
    \item El segmento anterior ya no existe (fue marcado para
        eliminación y todos los procesos hicieron \texttt{shmdt}).
    \item \texttt{shmget} con \texttt{IPC\_CREAT} crea uno
        \textbf{nuevo} con la misma clave, pero su contenido es
        \emph{indeterminado} (típicamente, todo ceros en Linux).
    \item \texttt{printf("Data read from memory: \%s\textbackslash{}n",
        str)} imprime una línea vacía.
\end{itemize}

\begin{lstlisting}[style=shellstyle]
$ ./shmreader     # primera vez: ve los datos
shared memory logical address: 0x7f...
Data read from memory: Hola desde shmwriter

$ ./shmreader     # segunda vez: segmento recien creado, vacio
shared memory logical address: 0x7f...
Data read from memory:
\end{lstlisting}

\subsection{9.b --- Comentar la destrucción y correr varias veces}

\begin{lstlisting}[style=cstyle]
// shmctl(shmid, IPC_RMID, NULL);   <-- comentada
\end{lstlisting}

Con esta modificación, el segmento \emph{persiste} entre ejecuciones.
Se mantiene activo en el kernel hasta que se borre explícitamente con
\texttt{ipcrm} o se reinicie el sistema:

\begin{lstlisting}[style=shellstyle]
$ ./shmwriter
Data to write to shmem: mensaje persistente
Data written in memory: mensaje persistente

$ ./shmreader
Data read from memory: mensaje persistente

$ ./shmreader     # de nuevo: los datos siguen ahi
Data read from memory: mensaje persistente

$ ./shmreader
Data read from memory: mensaje persistente

$ ipcs -m         # ver segmentos de memoria compartida activos
------ Memoria compartida ------
key        shmid   propietario  perms  bytes      nattch
0x4180...  98305   tomi         666    1024       0
\end{lstlisting}

Para limpiar manualmente:

\begin{lstlisting}[style=shellstyle]
$ ipcrm -m 98305
\end{lstlisting}

\begin{idea}
Esto ilustra una propiedad clave de la API: los segmentos IPC de
System V tienen \textbf{ciclo de vida del kernel}, no del proceso.
Sobreviven al cierre de todos los procesos que los usan, salvo que se
los elimine explícitamente. Es una diferencia importante con respecto
a \texttt{mmap MAP\_ANONYMOUS} o \texttt{shm\_open + mmap} (POSIX),
cuyo ciclo de vida se administra de manera más parecida a la de un
archivo.
\end{idea}

\subsection{9.c --- Lectura/escritura en \texttt{str + 500}}

Se trata de usar como base del dato compartido el offset $500$ dentro
del segmento. La modificación es simple y se aplica en ambos lados.

\paragraph{shmwriter (modificado).}

\begin{lstlisting}[style=cstyle]
char *str = (char *)shmat(shmid, (void *)0, 0);
printf("shared memory logical address: %p\n", str);
printf("Escribiendo en offset 500 (str+500 = %p):\n", str + 500);
printf("Data to write: ");
fgets(str + 500, 1024 - 500, stdin);
(str + 500)[strcspn(str + 500, "\n")] = '\0';
printf("Data written in memory at str+500: %s\n", str + 500);
shmdt(str);
\end{lstlisting}

\paragraph{shmreader (modificado).}

\begin{lstlisting}[style=cstyle]
char *str = (char *)shmat(shmid, (void *)0, 0);
printf("shared memory logical address: %p\n", str);
printf("Leyendo desde offset 500 (str+500 = %p):\n", str + 500);
printf("Data read from memory: %s\n", str + 500);
shmdt(str);
\end{lstlisting}

\paragraph{Notas.}

\begin{itemize}[leftmargin=*,nosep]
    \item \texttt{str + 500} es una dirección lógica válida porque el
        segmento tiene $1024$ bytes y $500 < 1024$.
    \item La aritmética de punteros funciona porque \texttt{str} es
        \texttt{char *}: \texttt{str + 500} avanza exactamente $500$
        bytes.
    \item En el caso del \texttt{writer} se acota
        \texttt{fgets(str + 500, 1024 - 500, stdin)} para que no se
        escriba más allá del fin del segmento.
    \item Los primeros $500$ bytes del segmento siguen sin tocar; si
        algún programa anterior los había escrito, su contenido
        permanece.
\end{itemize}

\subsection*{Resumen del ejercicio}

\begin{itemize}[leftmargin=*,nosep]
    \item La API \texttt{shmget}/\texttt{shmat}/\texttt{shmdt}/
        \texttt{shmctl} expone explícitamente la abstracción de
        \emph{región de memoria compartida} entre procesos.
    \item Cada proceso mapea la misma región a (posiblemente)
        \emph{distintas direcciones lógicas}, pero el contenido es
        compartido porque las \texttt{pte} apuntan a los mismos
        frames físicos.
    \item El ciclo de vida del segmento depende del sistema, no de
        los procesos: si nadie llama a \texttt{shmctl(IPC\_RMID)}, el
        segmento queda hasta el reboot.
    \item Es responsabilidad del programador sincronizar accesos al
        área compartida (semáforos, mutex en \emph{shm}, etc.).
\end{itemize}

% =====================================================================
\section{Conclusiones}
% =====================================================================

A lo largo del práctico se recorrió, ejercicio a ejercicio, toda la
maquinaria que el sistema operativo usa para gestionar la memoria:

\begin{itemize}[leftmargin=*]
    \item Los \textbf{Ejercicios 1 y 3} mostraron el \emph{lado
        práctico}: cómo medir el consumo real de memoria desde la
        línea de comandos y desde un programa propio, distinguiendo
        \texttt{VIRT}/\texttt{RES}/\texttt{PSS}/\texttt{USS} y los
        distintos segmentos del proceso (\texttt{.text},
        \texttt{.data}, heap, stack).
    \item Los \textbf{Ejercicios 2, 4, 5 y 6} se concentraron en la
        \emph{asignación} y el \emph{paginado}: el \emph{buddy
        system} ilustra los compromisos entre fragmentación interna y
        externa; el cálculo de la tabla de páginas en arquitecturas de
        $16$ y $32$ bits muestra por qué hacen falta tablas multinivel
        y/o \emph{huge pages} en sistemas reales.
    \item Los \textbf{Ejercicios 7, 8 y 9} aplicaron la maquinaria al
        problema de \emph{compartir memoria entre procesos} y mapear
        archivos en memoria. Los detalles relevantes (bit \texttt{D}
        para persistencia de \texttt{mmap}; ciclo de vida del
        segmento de \emph{shared memory}; visión común de distintos
        \texttt{vpn} sobre los mismos \texttt{ppn}) son consecuencia
        directa del modelo del capítulo teórico.
\end{itemize}

\begin{idea}
Lo que une todos los ejercicios es la \emph{misma} abstracción: la
tabla de páginas. Sirve para confinar un proceso, compartir memoria,
mapear archivos, hacer crecer el stack y abaratar \texttt{fork()}. La
diferencia entre cada caso está sólo en cómo el SO configura las
\texttt{pte} y reacciona ante \emph{page faults}, no en una
maquinaria nueva por feature.
\end{idea}

\end{document}
